\section{Rebuild from zero}\label{sec:rebuild}
Follow this section in order to go from nothing to the system as it was on 14 October 2026. Each step
ends with a check. Do not go on until it passes. The later sections give the details for each step.

\subsection{What ``the current version'' is}
\begin{table}[H]\centering\small
\begin{tabularx}{\linewidth}{lX}
\toprule
Part & Version \\
\midrule
Firmware & v1.0.0, protocol 1 (\cmd{bumpem info} prints \cmd{proto=1,fw=1.0.0}) \\
Server & 0.4.0 (\cmd{GET /api/v1/info} prints \cmd{"server": "0.4.0"}) \\
Repositories & \fillin{git tag of both repos, e.g.\ \cmd{v2026-10-14}} \\
ESCON configuration & \file{escon/moduleA_left_2026-10-01.edc} loaded and auto-tuned on every driver \\
Carrier board & v1, \file{hardware/carrier/fab/bumpem-carrier-gerbers.zip} \\
Printed parts & \file{hardware/enclosures/out/*.stl}, \cmd{KEY\_CLEAR} 0.2\,mm \\
\bottomrule
\end{tabularx}
\end{table}

\subsection{Step 1: get the files}
\begin{enumerate}
\item Clone both repositories into the \emph{same} folder (the server finds the UI at \file{../BumpemUI}):
\begin{lstlisting}
mkdir IBMS_Offenburg && cd IBMS_Offenburg
git clone https://github.com/lore95/BumpemServer.git
git clone <BumpemUI repository URL>
\end{lstlisting}
\fillin{BumpemUI repository URL (not on GitHub yet)}
\item The reference material is not in git (too large): \file{docs/thesis/}, \file{docs/reference/}
(datasheets), \file{docs/hardware-photos/}, \file{escon/*.pdf}, \file{legacy/}. Copy them into the same
paths. \fillin{where they are stored: shared drive / USB stick / ask whom}
\item Stanford's original files (CAD of motor mount and reel drum, full BOM, build guide):
\url{https://biomechatronics.stanford.edu/bump-em}.
\end{enumerate}
\textbf{Check:} \file{BumpemServer} and \file{BumpemUI} sit side by side; \file{docs/reference/} holds the
datasheets.

\subsection{Step 2: install the PC software}
\begin{table}[H]\centering\small
\begin{tabularx}{\linewidth}{lllX}
\toprule
Tool & Needed for & OS & Install \\
\midrule
Miniconda (Python 3.12) & server, UI, tests & all & \url{https://docs.conda.io/en/latest/miniconda.html} \\
PlatformIO & flashing the Teensy & all & \cmd{pip install platformio} (inside the env) \\
Teensy udev rule & flashing, serial port & Linux only & \file{00-teensy.rules} from pjrc.com \\
ESCON Studio & ESCON configuration & Windows only & maxon website (ESCON software package) \\
Chrome or Firefox & UI & all & \\
KiCad 8 & only to \emph{change} the carrier board & all & kicad.org \\
CadQuery, VTK & only to \emph{change} printed parts & all & \cmd{pip install cadquery matplotlib vtk} \\
\bottomrule
\end{tabularx}
\end{table}
\begin{lstlisting}
conda create -n BumpemHome python=3.12
conda activate BumpemHome          # every new terminal
cd BumpemServer
pip install -e ".[server,dev]"
pip install platformio
pytest                             # must print: 52 passed
bumpem serve --board sim           # then open http://127.0.0.1:8000/
\end{lstlisting}
\textbf{Check:} 52 tests pass; in the browser, \emph{Connect} to the simulator, \emph{Arm}, \emph{Perturb now}:
the live plot shows the pulse.

\subsection{Step 3: control board}
\begin{enumerate}
\item Have the carrier board made and soldered (\file{hardware/carrier/README.md}; Sect.~\ref{sec:carrier}).
Meter checks with the board bare: 3.3\,V $\leftrightarrow$ GND no beep; each plug pin to its Teensy socket pin
(EN 27--30, force 14--17) beeps; neighbouring pins do not.
\item Insert the Teensy 4.1 and the MCP4728, USB to the PC, nothing else connected.
\item Flash firmware v1 (erases what was on the Teensy):
\begin{lstlisting}
bumpem ports                       # find the Teensy, e.g. /dev/cu.usbmodem..., COM7
scripts/flash.sh                   # press the Teensy button if asked
bumpem info --port PORT            # proto=1, fw=1.0.0
bumpem monitor --port PORT         # about 100 lines per second
\end{lstlisting}
\item Optional desk test, no modules: jumper DAC VA--VD to pins 14--17, then
\cmd{bumpem set kp\_track 0 --port PORT}, the same for \cmd{kd\_track}, then \cmd{bumpem arm --port PORT}:
each channel reads about 31\,N (DAC 0.51\,V). Remove the jumpers afterwards.
\end{enumerate}
\textbf{Check:} \cmd{bumpem info} answers; telemetry at 100 lines/s. If it says ``DAC init failed'', check
DAC VCC, GND, SDA (18), SCL (19).

\subsection{Step 4: each tower}
Do one tower completely before starting the next. Nobody near the rope.
\begin{enumerate}
\item \textbf{Mechanics.} Motor, coupling, reel drum (radius 19.05\,mm), rope, eyebolt and load cell as in the
Stanford build guide pp.\,11--13. Mount the electronics with the printed parts (Sect.~\ref{sec:printed}).
\fillin{photos of one assembled IBMS tower: drum, rope path, load cell position, harness attachment}
\item \textbf{Power wiring} as in Fig.~\ref{fig:power}. Chopper DIP: S3 and S5 ON (55\,V). Mains side by the
electrician. Then the meter checks in Table~\ref{tab:powerchecks}.
\item \textbf{Motor} to ESCON J2/J3 with the colours of Sect.~\ref{sec:cabling}.
\item \textbf{ESCON configuration} (Sect.~\ref{sec:escon}): load module A's file, download, auto-tune, then
ESCON Studio $\rightarrow$ Diagnostics: all tests pass. Save as \file{escon/module<X>_<side>_<date>.edc}.
\item \textbf{Amplifier.} 24\,V to the IAA100, load cell to its sensor side. Set it so that 0--200\,N gives
0--3.3\,V, and never more than 3.3\,V. \fillin{IAA100 DIP switch positions read from a working amplifier;
load-cell wire colours}
\item \textbf{Signal cable}: the 5 wires (EN, SP+, SP$-$, GND, F) with ferrules into the module's MSTB plug,
plug into the carrier board.
\end{enumerate}

\begin{table}[H]\centering\small
\begin{tabularx}{\linewidth}{llX}
\toprule
\# & Meter & Pass \\
\midrule
T0 & DC V, supply +V to $-$V & about 0\,V (else wait for the capacitors) \\
T2 & continuity, supply +V to $-$V & no steady beep (a short chirp is the capacitors) \\
T3 & continuity & +V$\rightarrow$chopper 4, $-$V$\rightarrow$5, earth$\rightarrow$3, chopper 6$\rightarrow$J1+, 7$\rightarrow$J1$-$ all beep; +V to 5 does \emph{not} \\
T4 & continuity, +V to earth & no steady beep \\
T6 & look at the chopper DIP & S3 + S5 ON = 55\,V \\
T7 & power on via the switched strip & chopper LEDs off, ESCON LED on, chopper 4--5 reads 47--49\,V \\
\bottomrule
\end{tabularx}
\caption{Power checks before the first power-on, mains plug out for T0--T6 (full sheet:
\file{docs/tests/power-checks-module-C.md}).}
\label{tab:powerchecks}
\end{table}

\textbf{Check per tower:} ESCON LED slow green with enable low; Diagnostics passed.

\subsection{Step 5: acceptance (the system is at the current version when all pass)}
\begin{enumerate}
\item \cmd{bumpem serve --testing --board PORT --modules ABCD}, testing page, 48\,V on: Pull 5\,N then
10\,N, 400\,ms on each module in turn. Only that module's rope moves.
\item 24\,V on, ropes slack: each force reads a few N at rest; a hand pull changes only its own reading
(\file{scripts/force_check.py}).
\item Weights 6 / 8 / 12\,kg on each module: reading within about 2.5\,\% (Sect.~\ref{sec:calibration}).
\item Production server, \emph{Zero sensors}, \cmd{kff} 2, baseline 2\,N, \cmd{fmax\_n} 40, \emph{Save as
default}. Replug the Teensy, Connect: the values are back.
\item Arm one module at a time: it holds 2\,N. \emph{Record} a 10\,N perturbation: the CSV appears in
\file{Logs/}.
\item E-stop: press, every ESCON LED goes off; release, they come back.
\end{enumerate}
State on 14 October 2026: item 1 passed on A, C and D; A and D held 2\,N (item 5 without recording).
Items 2, 3, 4 and 6 and recording on the real board are still open; B waits for its ESCON.
